\begin{frame}{고차원으로 옮기면 보이는 것: 동심원}
  $x=(x_1,x_2)$ 를 $\phi(x) = (x_1, x_2, x_1^2+x_2^2)$ 로 올리면:
  \begin{itemize}
    \item $x_1^2+x_2^2$ = \textbf{원점으로부터의 거리 제곱}이므로
    \item 안쪽 원환체(반경 $\sim0.4$)의 모든 점 $\to$
      한 고도 $z \approx 0.4^2 = 0.16$ 에,
    \item 바깥쪽 원환체(반경 $\sim1.0$)의 모든 점 $\to$
      다른 고도 $z \approx 1.0^2 = 1.0$ 에 붙는다.
    \item 2차원에서 원 둘레를 따라 ``동일한'' 점이던 것이,
      세 번째 축 $z$를 추가하는 순간 두 개의 \textbf{수평 고도}로
      벌어진다 — \textbf{수평 평면} $z=0.5$ 하나로 완벽히 분리.
    \item 3차원의 평면은 원래 공간으로 투영하면 \textbf{원}이 되므로,
      ``3차원의 평평한 경계'' $\equiv$ ``2차원의 원형 경계''.
  \end{itemize}
  \vskip0.4em
  {\footnotesize \texttt{make\_circles}에선 바깥쪽 원환체=클래스 0,
  안쪽=클래스 1 — 라벨 번호는 임의, ``두 고도로 벌어진다''는
  구조는 동일.}
\end{frame}
